\subsection*{A shared layout recovers physical form}
A complete probability recount found that the earlier patch writer assigned
more mass to the blank than to the actual successor on 329 of its 348 blank
outputs. We therefore conditioned all patches on a common coarse image.
On the same 668 development contexts, known nonempty glyph matches rose from
198 to 649 and blanks fell from 348 to one (Fig.~\ref{fig:shared-layout}).
Exact matches to Han ideographs found in the source rose from 119 to 340.
Exact recorded successors remained at 18. The writer retained nine earlier
matches, gained nine and lost nine (paired difference 0.0 percentage points,
95\% article interval $-1.20$ to $+1.20$). Rotation of the training
successors produced 668 known nonempty glyphs and zero correct successors.
For the question about the history of the character for speech, the ordered
writer produced a letter followed by repeated ovals.

Training exposed 5,276 distinct target bitmaps and 539 coarse layouts.
We enumerated the greedy output for every selectable layout. Their range
contained the reference bitmap at only 242 of the 668 positions. Even perfect
choice of layout could not produce the other 426 reference images with this
writer. All four ordered completions reached the cap. A two-glyph context
also conflates different questions with identical visible endings. Preserving
finer distinctions in the writer and access to the earlier question in the
predictor are both necessary for this route to independent language.
